\documentclass[10pt,a4paper]{amsart}
\usepackage[margin=19mm]{geometry}
\usepackage{amsmath,amssymb,mathtools}
\usepackage[scr=rsfs,cal=euler]{mathalfa}
\usepackage{xfrac}
\usepackage[unicode,hidelinks,bookmarks=false]{hyperref}
\newcommand{\ie}{i.e.\ }
\newcommand{\cf}{cf.\ }
\newcommand{\resp}{respectively\ }
\newcommand{\Subsection}{\subsection}
\providecommand{\nobreakdash}{\nobreak\hbox{-}}
\setlength{\emergencystretch}{2em}
\multlinegap0pt
\usepackage{amssymb}
\usepackage{tikz}
\newtheorem{defn}{Definition}
\newtheorem{lem}{Lemma}
\title{Widths of Complements of Skeleta}
\author{Elliot Gathercole}
\date{Source: arXiv:2602.06863v1 (CC BY 4.0)}
\begin{document}
\maketitle

\noindent\emph{Source and licence.} Widths of Complements of Skeleta. Version arXiv:2602.06863v1 (CC BY 4.0), distributed by arXiv under the Creative Commons Attribution 4.0 International licence, http://creativecommons.org/licenses/by/4.0/, as recorded in the arXiv licence metadata for this exact version and shown on the abstract page https://arxiv.org/abs/2602.06863v1. Full licence notice: \texttt{licenses/arxiv-2602.06863-CC-BY-4.0.txt}.\par\smallskip

\noindent\textit{Editorial scope.} Counted unit: the article's lem integral\_orbit\_index (source0.tex lines 391--397). The article's complete original proof of it is delivered with it (source0.tex lines 398--419, one \texttt{proof} environment of 22 lines). No mathematics is written, repaired or re-derived: the statement and the proof are the article's own lines, in the article's own notation, and no retained line is substituted or re-flowed.  Retained uncounted as the source-local context the proof needs: lines 342--353.  Nothing was omitted from any retained span, so the package carries no omission record.  No retained line contains a citation, so the wrapper carries no bibliography and no bibliography entry.  No retained line contains a figure environment, an image-inclusion command or drawing code, so no image is delivered and the package declares no asset dependency.  Editorial formatting only, in addition: an AMS \texttt{amsart} preamble that restates the article's own declarations and macros used by the retained lines, namely the environments \texttt{defn}, \texttt{lem} on separate counters, together with the standard packages those lines need.  The retained environments print in this document as Definition 1, Lemma 1 in order of appearance, whereas the article prints the same items with the numbering of its own full document.  The complete native source of the article as distributed in the arXiv e-print for this exact version is bundled unchanged as \texttt{sources/194127/source0.tex}.
\begin{defn}[Outer Caps]
\label{outer_caps}
For any $S$-periodic Reeb orbit $\gamma$ in $(Y^R,\alpha)$, an outer cap is any representative of the class in $\delta\in\pi_2(M,\gamma)$ with boundary $[\gamma]$ defined as follows. We first define a rational class
\begin{equation}
    \delta=\sum_{v\in I(\gamma)}Sa_v(\gamma)\delta_v\in \pi_2(M,O(\gamma))\otimes\mathbb{Q},
\end{equation}
then show that $\delta\in\pi_2(M,O(\gamma))$.

By assumption on the diameter of the equivariant Darboux charts, we may assume that for all $v\in I(\gamma)$, we can represent $\delta_v$ by a disc contained in $B$, where $B$ is an embedded ball (we can choose $B$ to be a ball of diameter $<inj(M)$ where $inj(M)$ denotes the radius of injectivity of $M$ with respect to some metric). Because the boundary map $\pi_2(B,O(\gamma))\rightarrow \pi_1(O(\gamma))$ is an isomorphism, and the image of $\delta$ is the integral class $[\gamma]$ by construction, we have that $\delta\in\pi_2(M,O(\gamma))$ as required.

We denote by $u_{out}(\gamma)$ some representative of $\delta$ which bounds $\gamma$.
\end{defn}

\begin{lem}
\label{integral_orbit_index}
    For any $S$-periodic Reeb orbit $\gamma$ in $(Y^R,\alpha)$, if $Sa(\gamma)\in\mathbb{N}^{I(\gamma)}$, then
    \begin{equation}
        \mu_{CZ}^{\tau_{out}}(\gamma)+2S\sum_{v\in I(\gamma)}a_v(\gamma)w_v\in [-2n+2,2n-2].
    \end{equation}
\end{lem}

\begin{proof}
By hypothesis, $\gamma$ is a $1$-periodic orbit of the Hamiltonian $\mathbb{R}/\mathbb{Z}$-action generated by
\begin{equation}
G=S\sum_{v\in I(\gamma)}a_v(\gamma)r_v.
\end{equation}
Let $x=\gamma(0)$. Let $\tau_{eq}$ denote any trivialisation of $\gamma^*\xi$ arising from $D\phi^t_G$, i.e. $\tau_{eq}(t)=(D\phi_G^{S^{-1}t})_x|_\xi^{-1}:\xi_{\gamma(t)}\rightarrow \xi_x$, where we identify $\xi_x\cong\mathbb{C}^{n-1}$ in a fixed way.

First, we calculate $\mu_{CZ}^{\tau_{eq}}(\gamma)$. Let $A(t)=D(\phi^{St}\circ \phi_G^{-t})_x|_\xi$, so $\mu_{CZ}^{\tau_{eq}}(\gamma)=i(A)$, where we view $A$ as a map $[0,1]\rightarrow Sp(\xi_x)\cong Sp(2n-2)$.

The Lie algebra element $\dot{A}=\frac{dA}{dt}(0)$ is given by
\begin{equation}
\label{lie_element}
    \dot{A}=S\sum_{v\in I(\gamma)}\left(da_v\otimes X_{r_v}\right)_x|_\xi\in Sp(\xi_x)
\end{equation}
which satisfies $\dot{A}^2=0$ because for each $u,v\in I(\gamma)$, $\{a_u,r_v\}=0$, so $A(t)=id+t\dot{A}$. It follows from homotopy invariance that $i(A|_{[0,t]})$ is independent of $t\in (0,1]$. Letting $t\to 0$, we see that $i(A)\in [-2n+2,2n-2]$ because $A|_{[0,t]}$ becomes arbitrarily close to the constant path $id$.

Because $\mu_{CZ}^{\tau_{out}}(\gamma)=\mu_{CZ}^{\tau_{eq}}(\gamma)-2c_1^{\tau_{eq}}(u_{out}(\gamma))$, it suffices to show that $c_1^{\tau_{eq}}(u_{out}(\gamma))=S\sum_{v\in I(\gamma)}a_v(\gamma)w_v$.  First, we can extend $\tau_{eq}$ to a $\mathbb{R}^{I(\gamma)}/\mathbb{Z}^{I(\gamma)}$-invariant trivialisation of $\xi|_{O(\gamma)}$. The cap $u_{out}(\gamma)$ is homotopic in $O(\gamma)$ to a concatenation of discs $Sa_v(\gamma)u_v$ for $v\in I(\gamma)$ where $u_v$ is as in Definition \ref{outer_caps}, so
\begin{equation}
    c_1^{\tau_{eq}}(u_{out})=S\sum_{v\in I(\gamma)}a_v(\gamma)c_1^{\tau_{eq}}(u_v).
\end{equation}
The trivialisation $\tau_{eq}$ is equivalent as a trivialisation of $\gamma_v^*\xi$, where $\gamma_v=\partial u_v$, to that induced by an $r_v$-equivariant Darboux chart, so $c_1^{\tau_{eq}}(u_v)=w_v$, as required.
\end{proof}

\end{document}
